\bigskip\Needspace{13\baselineskip}
\pdfbookmark[1]{Tafeln für die Störungen der Ceres. März 1803}{ceresstoerungen180303}
\begin{center}
TAFELN FÜR DIE STÖRUNGEN DER CERES.\par
\medskip\hrule\medskip
Monatliche Correspondenz. Band VII, S.\ 259—262. März 1803.\par
\medskip\hrule\bigskip
GAUSS AN VON ZACH.
\end{center}\medskip
\begin{flushright}Braunschweig, 26. Dezember 1802.\end{flushright}

Vor einiger Zeit schrieb ich Ihnen, dass ich darüber nachgedacht habe, die Tafeln für die Störungen der \textit{Ceres} zu kürzen, und Ihnen später eine Probe davon vorlegen wolle. (Mon. Corr. Band VI, S.\ 498) [S.\ 230 dieses Bandes] Ich habe jetzt die Ehre, dies näher zu erläutern und Ihnen statt einer Probe eine Tafel \textit{in extenso} zu übersenden.

Schon jetzt, da ich mich noch auf die Störungen durch Jupiter und auf die ersten Potenzen der Exzentrizitäten beschränkt habe, gibt es so viele Gleichungen, dass es äußerst mühsam wäre, ihre Zahlenwerte für eine große Zahl von Positionen ohne Tafeln aus den Formeln selbst zu berechnen. Aber auch mit Tafeln herkömmlicher Anordnung ist die Arbeit wirklich mühsam. Für \textit{alle} Gleichungen, die im Novemberheft 1802 der M.\ C. [S.\ 228 dieses Bandes] abgedruckt sind, wären \textit{vierzig} Tafeln nötig, und noch immer fast \textit{dreißig}, wenn man die kleineren Gleichungen weglässt, die einzeln weniger als $2''$ betragen. So viele Argumente zu bilden, in so vielen Tafeln nachzuschlagen und so viele Additionen auszuführen, verlangt immer noch viel Geduld, wenn man es oft tun muss.
Ich habe daher versucht, diese Tafeln auf eine handlichere Weise zu kürzen. Die \textit{acht} ersten Gleichungen \textit{für die Länge} (wobei nebenbei gesagt der erste Koeffizient auf S.\ 495 [228] statt $230''00$ vielmehr $231''00$ lauten muss) ergeben von selbst \textit{eine} Tafel mit dem Argument $\CeresSymbol-\JupiterSymbol$. Die übrigen vierzehn würden aber ebenso viele Tafeln erfordern. Man könnte sie in einer Tafel mit zwei Eingängen unterbringen. Sie wäre jedoch so umfangreich, so mühsam zu berechnen und so umständlich zu benutzen, dass daran nicht zu denken ist. Ich habe daher folgendes Verfahren gewählt:
\SourceLeaf{241}%

Bezeichne ich der Kürze halber $\CeresSymbol-\JupiterSymbol$ mit $D$, so sieht man bei einiger Aufmerksamkeit, dass sich jede dieser vierzehn Gleichungen in die Form
\[
a\sin(iD-\CeresSymbol-A)=a\sin(iD-A)\cos\CeresSymbol-a\cos(iD-A)\sin\CeresSymbol
\]
bringen lässt, wobei $i$ positive und negative ganze Zahlen bezeichnet. Die Summe aller Gleichungen wird also durch $P\cos\CeresSymbol-Q\sin\CeresSymbol$ ausgedrückt, wobei $P$ und $Q$ Funktionen von $D$ sind. Setzt man
\[
\frac PQ=\operatorname{tang}\varphi,\quad
\frac P{\sin\varphi}=\frac Q{\cos\varphi}=R,
\]
so hängen auch $R$ und $\varphi$ nur von $D$ ab, und die Summe aller vierzehn Gleichungen
\[
=R\sin(\varphi-\CeresSymbol)
\]
Man kann also in einer Tafel mit dem Argument $D$ für alle Werte von $D$ die entsprechenden Werte von $R$ und $\varphi$ eintragen und erhält durch eine sehr leichte Rechnung sofort die Summe all dieser Gleichungen. Ich habe es jedoch vorgezogen, auf die Angabe von $\varphi$ zu verzichten und dafür $\varphi-2D$ einzutragen, weil sich dieser Winkel weniger stark ändert als $\varphi$ und nur periodisch innerhalb gewisser Grenzen schwankt, während $D$ alle Werte von $0$ bis $360^0$ durchläuft. Dagegen würde $\varphi$ in dieser Zeit zwei Umläufe vollziehen. Dieses $\varphi-2D$ wird in der dritten Tafel mit $B$ bezeichnet, ebenso wie das obige $R$ mit $A$. Auf diese Weise ist also die Summe aller von der einfachen Exzentrizität abhängigen Gleichungen für die Länge
\[
A\sin(B+2D-\CeresSymbol)=A\sin(B+\CeresSymbol-2\JupiterSymbol)
\]

Entsprechend stehen die von der Exzentrizität unabhängigen Gleichungen des Radiusvektors in Tafel II, und die von den einfachen Exzentrizitäten abhängigen sind in Tafel IV in die Form $A\sin(B+2\CeresSymbol-3\JupiterSymbol)$ gebracht. In derselben Form stehen schließlich in Tafel V alle im Novemberheft [S.\ 229 dieses Bandes] aufgeführten Gleichungen für die Breite.

Auf diese Weise haben wir also nur fünf Tafeln. Dabei ist zwar etwas trigonometrische Rechnung nötig, die mir aber jedenfalls viel leichter fällt, als wenn ich jede der hier zusammengefassten Gleichungen einer eigenen Tafel entnehmen müsste. Auch ließen sich noch einige vorteilhafte Änderungen vornehmen, wenn man z.\ B. statt $A$ gleich $\log A$ eintrüge, \textit{oder auch auf andere Weise}. Das will ich mir aber für später aufheben, denn diese Tafeln sollen überhaupt nur so lange dienen, bis die Elemente anhand künftiger Beobachtungen weiter berichtigt und die Störungen danach genauer und vollständiger berechnet sind. Vorerst denke ich, dass mir diese Tafeln besonders bei der Berechnung der Beobachtungen des nächsten Jahres gute Dienste leisten sollen.

\Needspace{7\baselineskip}
\begin{center}\textit{Anmerkungen.}\end{center}
1.\quad $\CeresSymbol$ und $\JupiterSymbol$ bezeichnen die mittleren heliozentrischen Längen der Ceres und des Jupiter.

2.\quad Die Störungen des Radiusvektors sind in Millionsteln des Radius der Erdbahn ausgedrückt.

3.\quad Bei den Störungen der Breite und den zweiten Teilen der Störungen der Länge und des Radiusvektors ist $A$ als positiv anzusehen. Beim Sinus von $B+\CeresSymbol-2\JupiterSymbol$ und von $B+2\CeresSymbol-3\JupiterSymbol$ ist das Vorzeichen entsprechend zu beachten.

4.\quad Für die Breite nimmt man $B$ bei den ersten sechs Zeichen des Arguments oder beim Absteigen aus der ersten Spalte, bei den letzten sechs Zeichen oder beim Aufsteigen dagegen aus der dritten Spalte. $A$ ist beiden gemeinsam. Das geht schon aus der Anordnung der Tafel hervor.
